\subsection{Two public EEG acquisitions as negative exact-epoch integrity controls}
\label{sec:v28-negative-acquisition-controls}
The exact-epoch cross-identifier finding is not presented
as a universal property of public EEG repositories.
We preselected two previously exposed independent
recording acquisitions (not new performance-validation
cohorts) before conducting the content identity scan.
Weibo2014 was processed with its original fixed
C3/Cz/C4 three-channel, left/right motor imagery
pipeline. BNCI2014\_002 retained the original native
EEG channels and fixed right-hand/feet task. Both
used the original MOABB 1.7.2, 8--30 Hz,
160 Hz, source-compatible preprocessing.

\begin{table}[tb]
\centering\small
\caption{Exact processed full EEG trial-content identity across
all preselected subject identifiers in three public acquisitions.
Only exact repeats under these fixed pipelines are audited;
no new AUC is evaluated in this table.}
\label{tab:v28-integrity-controls}
\begin{tabular}{lrrrr}
\toprule
Acquisition & IDs & Trials & Unique trials & Cross-ID edges\\
\midrule
Liu2024 & 50 & 2,000 & 1,960 & 1\\
Weibo2014 & 10 & 1,580 & 1,580 & 0\\
BNCI2014\_002 & 14 & 2,240 & 2,240 & 0\\
\bottomrule
\end{tabular}
\end{table}

The two negative-control audits completed successfully in
GitHub Actions run \texttt{38289332514},
receipt SHA-256 prefix \texttt{eb870591f13e7593}
(full result and digest in its immutable artifact).
For both controls all trial arrays were unique under
the exact numeric and event-label equality criterion.
This constrains the observed duplicated released
record to the specific Liu2024 acquisition and
preprocessing representation. It does not establish
that other releases are free of exact duplicates,
nor do negative exact-equality tests exclude near-duplicates
or nonidentical overlapping signal windows.
The two control cohorts had already been used
in earlier DefaultReg method development,
so they do not add fresh independent algorithm
generalisation evidence.
